% !TEX root = CSL-main.tex
\section{Cyclic Proofs and the Global Trace Condition}
\label{sec:cyclic}

In this section we define the global trace condition (GTC), a validity criteria used to verify the soundness of circular proofs. As our primary examples, we consider the cyclic proof systems $\CLKIDomega$ and $\CLJIDomega$, and show how the GTC for both systems is equivalent to the SCT condition in Definition \ref{def:SCT}. Our approach is motivated by Ikebuchi's work \cite{Mirai}, where she formally stated the equivalence between the GTC and SCT, by extracting the size-change graphs of a cyclic proof. In our case, we extract a transition graph with bi-partite relations of a cyclic proof, which in principle corresponds to the same process. We begin by introducing the notion of cyclic pre-proofs.

\subsection{Cyclic Pre-Proofs}
\label{subsec:cyclicPreProofs}

$\CLKIDomega$ and $\CLJIDomega$ are both first order logics, hence we will consider a first-order signature $\Sigma$, which contain denumerable many constant, function and predicate symbols. Terms, formulas and $\Sigma$-structures are defined as usually for first order logic. We consider sequents as originally defined by Gentzen using finite lists. In the case of $\CLKIDomega$, sequents have the form $\Gamma \vdash \Delta$, while in $\CLJIDomega$, they are defined as $\Gamma \vdash \phi$, where $\Gamma$ and $\Delta$ are finite lists of formulas, and $\phi$ is list containing the single formula $\phi$. Moreover, the comma between lists of formulas denotes list concatenation (e.g. $\Gamma, \Delta$). For any formula $F$ in a list $\Gamma$, we define the \emph{occurrence} of $F$ in $\Gamma$ as the formula $F$ paired with the position $i$ where it appears in $\Gamma$, written as $(F,i)$. Predicate symbols are divided between finitely many \emph{inductive predicates}, as in Martin-Löfs \emph{productions} \cite{MartinLof}, and denumerable many \emph{ordinary predicates}, defined as usual predicates. To formally define inductive predicates, we introduce the notion of \emph{inductive definition set}:

\begin{definition}[Inductive definition set]
    \label{defIndSet}
    An inductive definition set $\Phi$ is a finite set of \emph{productions} of the form:

    \begin{center}
        \centering
        \begin{prooftree}
            \hypo{Q_1\mathbf{u_1} \ldots Q_h\mathbf{u_h}}
            \hypo{P_1\mathbf{t_1} \ldots P_m\mathbf{t_m}}
            \infer2[]{P\mathbf{t}}
        \end{prooftree}
    \end{center}

    Where $Q_1, \ldots, Q_h$ are ordinary predicates symbols, $P_1, \ldots, P_m$ are inductive predicate symbols, and $\mathbf{u_1}, \ldots, \mathbf{u_h}$, $\mathbf{t_1}, \ldots, \mathbf{t_m}$, $\mathbf{u_1}, \ldots, \mathbf{t}$ are tuples of terms.
\end{definition}

\begin{figure}
    \centering
    \includegraphics[scale = 0.5]{example-image}
    \caption{Rules of LK.}
    \label{LKrules}
\end{figure}


\begin{example}
We define the inductive definition set $\Phi_{NEO}$, using the following productions rules:

\begin{center}
    \centering
    \begin{minipage}{0.10\textwidth}
        \centering
        \begin{prooftree}
            \hypo{}
            \infer1[]{N0}
        \end{prooftree}
    \end{minipage}
    \begin{minipage}{0.10\textwidth}
        \centering
        \begin{prooftree}
            \hypo{Nx}
            \infer1[]{Nsx}
        \end{prooftree}
    \end{minipage}
    \begin{minipage}{0.10\textwidth}
        \centering
        \begin{prooftree}
            \hypo{}
            \infer1[]{E0}
        \end{prooftree}
    \end{minipage}
    \begin{minipage}{0.15\textwidth}
        \centering
        \begin{prooftree}
            \hypo{Ox}
            \infer1[]{Esx}
        \end{prooftree}
    \end{minipage}
    \begin{minipage}{0.12\textwidth}
        \centering
        \begin{prooftree}
            \hypo{Ex}
            \infer1[]{Osx}
        \end{prooftree}
    \end{minipage}
\end{center}

Where $0$ is a constant representing the number zero and $s$ is a unary function representing the successor function. The predicate $N$ is used to inductively represent the natural numbers, while $E$ and $O$ are \emph{mutually defined} inductive predicates defining even and odd numbers respectively.
\end{example}

We define the proof system for $\CLKIDomega$ as the extension of the classical logic LK, defined in Figure \ref{LKrules}, with additional rules used to handle inductive predicates on the right and left part of the sequent. These rules for inductive predicates are defined as follows: let $P$ be an inductive predicate, and let $\Phi = \{\Phi_1, \ldots, \Phi_n\}$ be the inductive definition set, such that all $\Phi_i$ are defined as in Definition \ref{defIndSet} and have as conclusion $P$.

\begin{itemize}
    \item For each production $\Phi_i$, the system $\CLKIDomega$ includes a corresponding right introduction rule, denoted ($P\text{R}_i$):
    \begin{equation}
    \label{def:caseRuleRight}
    \begin{prooftree}
        \hypo{\Gamma \vdash Q_1\mathbf{u_1}, \Delta \ldots \Gamma \vdash Q_h\mathbf{u_h}, \Delta}
        \hypo{\Gamma \vdash P_1\mathbf{t_1}, \Delta \ldots \Gamma \vdash P_m\mathbf{t_m}, \Delta}
        \infer2[($P\text{R}_i$)]{\Gamma \vdash P\mathbf{t}, \Delta}
    \end{prooftree}
    \end{equation}


    \item There is a left introduction rule associated to $P$ in $\CLKIDomega$, denoted \emph{case split rule } or (Case $P$):
    \begin{equation}
    \label{def:caseRuleLeft}
    \begin{prooftree}
        \hypo{ \{ \Gamma, \mathbf{u} = \mathbf{t^i}, Q^i_1\mathbf{u^i_1}, \ldots , Q^i_h\mathbf{u^i_h}, P^i_1\mathbf{t^i_1}, \ldots, P^i_m\mathbf{t^i_m} \vdash \Delta \}_{1 \leq i \leq n} }
        \infer1[(Case $P$)]{\Gamma, P\mathbf{u} \vdash \Delta}
    \end{prooftree}    
    \end{equation}
    Where for each production $\Phi_i$, there is a corresponding sequent in the premises of the form $\Gamma, \mathbf{u} = \mathbf{t^i}, Q^i_1\mathbf{u^i_1}, \ldots , Q^i_h\mathbf{u^i_h}, P^i_1\mathbf{t^i_1}, \ldots, P^i_m\mathbf{t^i_m} \vdash \Delta$ subject to the restriction that any free variable in $\mathbf{u^i_1},...,\mathbf{u^i_h},\mathbf{t^i_1},...,\mathbf{t^i_m}$, or $\mathbf{t^i}$, does not occur in $\Gamma, \Delta$ and $\mathbf{u}$.
\end{itemize}

\begin{example}
For the previously defined inductive predicates $N$, $E$ and $O$, we can associate the following rules:    

\begin{center}
    \centering
    \begin{minipage}{0.2\textwidth}
        \adjustbox{scale=0.8}{
        \begin{prooftree}
            \hypo{}
            \infer1[($E\text{R}_1$)]{\Gamma \vdash E0, \Delta}
        \end{prooftree}
        }
    \end{minipage}
    \hfill
    \begin{minipage}{0.2\textwidth}
        \adjustbox{scale=0.8}{
        \begin{prooftree}
            \hypo{\Gamma \vdash Ox, \Delta}
            \infer1[($E\text{R}_2$)]{\Gamma \vdash Esx, \Delta}
        \end{prooftree}
        }
    \end{minipage}
    \hfill
    \begin{minipage}{0.2\textwidth}
        \adjustbox{scale=0.8}{
        \begin{prooftree}
            \hypo{\Gamma \vdash Ex, \Delta}
            \infer1[($O\text{R}_1$)]{\Gamma \vdash Osx, \Delta}
        \end{prooftree}
        }
    \end{minipage}
    \hfill
    \begin{minipage}{0.35\textwidth}
        \adjustbox{scale=0.8}{
        \begin{prooftree}
            \hypo{\Gamma, x = 0 \vdash \Delta}
            \hypo{\Gamma, x = sz, Nz \vdash \Delta}
            \infer2[(Case $N$)]{\Gamma, Nx \vdash \Delta}
        \end{prooftree}
        }
    \end{minipage}

\end{center}
\end{example}

For each inference rule in $\CLKIDomega$, we define the \emph{sub-occurrence} relation, which maps occurrences in the conclusion of the rule with occurrences in its premises. To simplify the presentation, we define this relation pictographically (inspired by the work in \cite{SaurinBauer}), to illustrate how occurrences are tracked through each inference.

\begin{definition}[sub-occurrence relation]
    Let $r$ be a rule in $\CLKIDomega$, with conclusion $\Gamma \vdash \Delta$ and premises $\Gamma_1 \vdash \Delta_1, \ldots, \Gamma_n \vdash \Delta_n$. Moreover, consider the ocurrences $(F,i)$ in $\Gamma \vdash \Delta$ and $(G,j)$ in $\Gamma_k \vdash \Delta_k$ for some $ 1 \leq k \leq n$. We say that $(G,j)$ is a sub-occurrence of $(F,i)$ in $r$, noted $(F,i) \rightarrow_{r} (G,j)$ if and only if there is an arrow drawn from $(F,i)$ to $(G,j)$ in the definition of rule $r$ in the Figure \ref{LKrules} or the Formulas \ref{def:caseRuleRight} and \ref{def:caseRuleLeft}. We say that $(F,i)$ is an active occurrence of the rule, if the arrow coming out from it, is colored in red.
\end{definition}

The proof system for $\CLJIDomega$ and its sub-occurrence relation, is naturally obtained by restricting the rules for inductive predicates over intuitionistic sequents, preserving the exchange rule, and replacing LK with its intuitionistic version, named LJ, which is presented in \cite{Berardi}. \emph{Cyclic pre-proofs} are used to formally represent the graph structure of a derivation in $\CLKIDomega$ and $\CLJIDomega$.

\begin{definition}[Cyclic pre-proof] A $\CLKIDomega$ (resp. $\CLJIDomega$) pre-proof of sequent $\Gamma \vdash \Delta$ (resp. $\Gamma \vdash \phi$), is a tuple $\Pi = (\mathcal{D}, \mathcal{F})$, where:

    \begin{itemize}
        \item $\mathcal{D} = (V, s, r, p)$ is a derivation tree where $V$ is a finite set of vertices, $s: V \rightarrow Seqs$ is a total function labelling vertices with sequents in $\CLKIDomega$ (resp. $\CLJIDomega$), $r: V \rightarrow Rules$ is a partial function labelling vertices with inference rules in $\CLKIDomega$ (resp. $\CLJIDomega$), $p: \mathbb{N} \times V \rightarrow V$ is a partial function that maps a number $n$ and a vertex $v$, with a vertex $v'$, such that $v'$ is labeled by the $n^{\text{th}}$ premise of the rule $r(v)$ over the sequent $s(v)$, and for every $v \in V$, 
            \begin{prooftree}
                \hypo{\vdash s(p_1(v)) \quad \dots \quad \vdash s(p_m(v))}
                \infer1[$(r(v))$]{\vdash s(v)}
            \end{prooftree} 
        is an instance of the rule $r(v)$ in $\CLKIDomega$ (resp. $\CLJIDomega$), with $m$ premises. A vertex $B \in V$ such that $r(B)$ is undefined is called a bud node. The set of bud nodes of a derivation tree $\mathcal{D}$, is written as $Bud(\mathcal{D})$.

        \item $\mathcal{F}: V \rightarrow V$ is a partial function, that maps each one of the bud nodes in $Bud(\mathcal{D})$ to an internal node in the derivation tree, called a companion, satisfying that if $B \in Bud(\mathcal{D})$ and $\mathcal{F}(B) = C$, then $r(C)$ is defined and $s(B) = s(C)$.
    \end{itemize}

Moreover, the graph of $\Pi$, written as $G_{\Pi}$, is the graph obtained after adding the edges induced by $\mathcal{F}$ and deleting the bud vertices. Specifically, $G_{\Pi} = (V', s, r, p')$ where, $V' = V \setminus Bud(\mathcal{D})$ and $p'$ is defined for any $j$ as: $p'_j(v) = \mathcal{F}(p_j(v))$ if $p_j(v) \in Bud(\mathcal{D})$, or $p'_j(v) = p_j(v)$ otherwise.
\end{definition}

\begin{example}
\begin{figure}
\centering

\adjustbox{scale=0.8}{
\begin{prooftree}
    \hypo{}
    \infer1[($E\text{R}_1$)]{\vdash E0, O0}

    \hypo{Nx \vdash Ox, Ex \ (\dag)}
    \infer1[(Subst)]{Ny \vdash Oy, Ey}
    \infer1[($O\text{R}_1$)]{Ny \vdash Oy, Osy}
    \infer1[($E\text{R}_2$)]{Ny \vdash Esy, Osy}
    \infer1[($=$L)]{x = sy, Ny \vdash Ex, Ox}
    \infer2[(Case $N$)]{Nx \vdash Ex, Ox \ (\dag)}
    \infer1[($\lor$R)]{Nx \vdash Ex \lor Ox}
\end{prooftree}
}

\caption{$\CLKIDomega$ cyclic pre-proof of the sequent $Nx \vdash Ex \lor Ox$. The function $\mathcal{F}$ is represented by labeling each bud and its corresponding companion node with the $(\dag)$ symbol.}
\label{exampleCircProof}
\end{figure}

Figure \ref{exampleCircProof} is an example of a $\CLKIDomega$ cyclic pre-proof, taken from \cite{BrotherstonSimpson}.
\end{example}

As mentioned before, the GTC is used as a criterion to verify the soundness of cyclic pre-proofs in $\CLKIDomega$ and $\CLJIDomega$. To define it, it is first necessary to formalize both the notions of \emph{trace} and \emph{progress point}. 

\begin{definition}[Trace and Progress Point]
    \label{def:traceCyclic}
    Let $\Pi$ be a $\CLKIDomega$ (resp. $\CLJIDomega$) pre-proof and let $(v_i)$ be a sequence of vertices defining a path in $G_{\Pi}$. A trace over $(v_i)$ is a sequence of occurrences $\uptau$ such that, for all $i$:

    \begin{itemize}
          \item $\uptau(i) = (P\mathbf{u},j)$ is in $\Gamma$, where $P$ is an inductive predicate, and $s(v_i) = \Gamma \vdash \Delta$ (resp. $s(v_i) = \Gamma \vdash \phi$).

          \item If $\uptau(i) = (P\mathbf{u},j)$, then $\uptau(i+1) = (Q\mathbf{t},k)$, where $(Q\mathbf{t},k)$ is an occurrence such that $Q$ is an inductive predicate, $(P\mathbf{u},j) \rightarrow_{r(v_i)} (Q\mathbf{t},k)$, and $(Q\mathbf{t},k)$ is in $\Gamma'$, where $s(v_{i+1}) = \Gamma' \vdash \Delta'$ (resp. $s(v_{i+1}) = \Gamma' \vdash \phi'$).

          \item If $r(v_i) = (\text{Case } P)$ is a case split rule, and $\uptau(i) = (P\mathbf{u},j)$ is the active occurrence of the rule, then $i$ is said to be a progress point of the trace.
    \end{itemize}  

\end{definition}

\begin{remark}
    Explain here: 1) our definition is different 2) however is in principle the definition of brotherston extended to ocurrences 3)
\end{remark}  



